% Authorial checks of the declared models; not transcriptions of medieval proofs.
\section{Checks for the declared models of supposition}
\label{sec:annex-p13-supposition}

These proofs concern the extension clauses and relational comparisons
declared in \cref{sec:p13-medieval-supposition-theories}. Their historical
starting points are Parsons's manuscript discussions of opposition and
conversion, the twenty-four moods, dependent terms, and infinitizing
negation \cite[pp.~161--168,176--179,183]{Parsons2007FordhamSupposition}.
They do not turn his composite historical account into one medieval
formal system. Throughout the first two proofs, \(S,M,P\) are subsets of
a fixed domain, with no assumption that any particular subset is nonempty.

\needspace{8\baselineskip}
\begin{annexitem}{annex:p13.supposition.square}{thm:p13-supposition-square}{The square and its conversions}
\begin{proof}
The definitions give \(O(S,P)=\neg A(S,P)\) and
\(E(S,P)=\neg I(S,P)\), so each indicated pair consists of
contradictories. If \(A(S,P)\), choose \(s\in S\), possible because
\(S\ne\varnothing\). Since \(S\subseteq P\), this element belongs
to \(S\cap P\); hence \(I(S,P)\). Contraposition of this implication
gives \(E(S,P)\Rightarrow O(S,P)\). Thus A and E cannot both hold.
If I and O were both false, E and A would both hold, the same contradiction.

Since \(S\cap P=P\cap S\), I and E are each equivalent to their
conversions. From A, the nonempty intersection just obtained also gives
\(I(P,S)\). From E, conversion gives \(E(P,S)\), and subalternation
then gives \(O(P,S)\).

For failure of O-conversion take distinct \(a,b\), put
\(S=\{a,b\}\) and \(P=\{a\}\). Then \(A(S,P)\) is false
because \(b\notin P\), so \(O(S,P)\) is true. But
\(A(P,S)\) is true, so \(O(P,S)\) is false. Finally, when
\(S=\varnothing\), its A-clause is false whatever \(P\) is;
therefore O is true without a subject witness.
\end{proof}
\end{annexitem}

\needspace{8\baselineskip}
\begin{annexitem}{annex:p13.supposition.moods}{thm:p13-supposition-moods}{All twenty-four displayed moods}
\begin{proof}
We verify the nineteen named rows in the order of the body table. The
five unnamed rows then follow by the square. In a negative conclusion,
exhibiting \(s\in S\setminus P\) suffices for \(O(S,P)\), because it
refutes \(S\subseteq P\). This is a sufficient argument when a witness
is available; it is not the definition of O in the empty-subject case.

\begin{enumerate}[label=\arabic*.]
\item \emph{Barbara.} The premises give \(S\ne\varnothing\),
  \(S\subseteq M\) and \(M\subseteq P\); hence
  \(S\ne\varnothing\) and \(S\subseteq P\), which is \(A(S,P)\).
\item \emph{Celarent.} If \(s\in S\cap P\), then
  \(S\subseteq M\) puts \(s\) in \(M\cap P\), contrary to
  \(E(M,P)\). Thus \(S\cap P=\varnothing\).
\item \emph{Darii.} A witness \(s\in S\cap M\) supplied by
  \(I(S,M)\) belongs to \(P\), because \(M\subseteq P\).
  It witnesses \(I(S,P)\).
\item \emph{Ferio.} Choose \(s\in S\cap M\). Since
  \(M\cap P=\varnothing\), this is a member of \(S\setminus P\).
\item \emph{Baralipton.} Choose \(p\in P\), using \(A(P,M)\).
  The inclusions \(P\subseteq M\subseteq S\) put this witness in
  \(S\cap P\).
\item \emph{Celantes.} From \(P\subseteq M\) and
  \(M\cap S=\varnothing\), it follows that \(P\cap S=\varnothing\),
  equivalently \(E(S,P)\).
\item \emph{Dabitis.} A witness in \(P\cap M\) also belongs to
  \(S\), since \(M\subseteq S\); it witnesses \(I(S,P)\).
\item \emph{Fapesmo.} The affirmative premise gives
  \(M\ne\varnothing\) and \(M\subseteq S\). Any \(m\in M\)
  is outside \(P\), by \(E(P,M)\); hence \(m\in S\setminus P\).
\item \emph{Frisesomorum.} A witness \(m\in M\cap S\) is outside
  \(P\), by \(E(P,M)\), and therefore belongs to \(S\setminus P\).
\item \emph{Cesare.} Since \(S\subseteq M\) and
  \(P\cap M=\varnothing\), we have \(S\cap P=\varnothing\).
\item \emph{Camestres.} Since \(P\subseteq M\) and
  \(S\cap M=\varnothing\), again \(S\cap P=\varnothing\).
\item \emph{Festino.} A witness \(s\in S\cap M\) cannot be in
  \(P\), because \(P\cap M=\varnothing\).
\item \emph{Barocho.} Suppose its conclusion \(O(S,P)\) were false.
  Then \(A(S,P)\). Together with \(A(P,M)\), Barbara gives
  \(A(S,M)\), contradicting the second premise \(O(S,M)\).
  This argument covers \(S=\varnothing\) too: in that case the
  supposed A-conclusion is already false by definition.
\item \emph{Darapti.} The two A-premises give a nonempty \(M\)
  contained in both \(S\) and \(P\). An element of \(M\) witnesses
  \(I(S,P)\).
\item \emph{Felapto.} The premise \(A(M,S)\) supplies an element of
  \(M\subseteq S\). The premise \(E(M,P)\) keeps it outside \(P\).
\item \emph{Disamis.} A witness in \(M\cap P\) is also in \(S\),
  since \(M\subseteq S\), and so witnesses \(I(S,P)\).
\item \emph{Datisi.} A witness in \(M\cap S\) is also in \(P\),
  since \(M\subseteq P\), and so witnesses \(I(S,P)\).
\item \emph{Bocardo.} From \(A(M,S)\), \(M\ne\varnothing\) and
  \(M\subseteq S\). Since \(O(M,P)\) denies
  \(M\ne\varnothing\land M\subseteq P\), it follows here that
  \(M\nsubseteq P\). Choose \(m\in M\setminus P\); then
  \(m\in S\setminus P\).
\item \emph{Ferison.} A witness in \(M\cap S\) lies outside
  \(P\), because \(M\cap P=\varnothing\).
\end{enumerate}

The unnamed row after Barbara replaces its A-conclusion by I, valid by
\cref{thm:p13-supposition-square}. Each unnamed row following Celarent,
Celantes, Cesare or Camestres replaces the proved E-conclusion by O,
also valid by that theorem. These are exactly the other five displayed
rows. Every witness used above comes from an I-premise or the explicit
nonemptiness conjunct of an A-premise. No inference extracts a witness
from an O-premise without first establishing its subject's nonemptiness.
\end{proof}
\end{annexitem}

\needspace{8\baselineskip}
\begin{annexitem}{annex:p13.supposition.scope}{thm:p13-supposition-scope}{Three finite relational countermodels}
\begin{proof}
For the first comparison let the people be \(a_1,a_2\) and the
donkeys be \(d_1,d_2\), all distinct. Interpret seeing, between these
ranges, by just \((a_1,d_1)\) and \((a_2,d_2)\). Each person sees a
donkey. No donkey is seen by both people: \(a_2\) does not see
\(d_1\), and \(a_1\) does not see \(d_2\). Thus
\[
 \forall a\in\{a_1,a_2\}\ \exists d\in\{d_1,d_2\}\ R(a,d)
\]
is true while the formula with the two quantifier blocks interchanged
is false. The common noun for a seeing-thing must retain its dependence
on the donkey argument; replacing it by a fixed unary range would lose
precisely the information that distinguishes the two rows of the relation.

For the second comparison use one girl \(g\), two presents
\(p_1,p_2\), and two boys \(b_1,b_2\). Let the giving relation
contain exactly \((g,p_1,b_1)\) and \((g,p_2,b_2)\), where the three
positions are giver, present and recipient. For each present the girl
has a recipient, so \(\forall p\,\exists b\,G(g,p,b)\) holds over
these ranges. Neither boy receives both presents, so
\(\exists b\,\forall p\,G(g,p,b)\) fails. Since there is just one
girl, adding an outer existential quantifier for the girl does not
alter this contrast.

Finally take distinct \(o,d,c\). Let \(o\) own \(d\), and let
the father-of relation contain \((d,c)\) but not \((d,o)\).
Then \(d\) is owned by \(o\) and is a father, because it is the
father of \(c\); it is not the father of \(o\). The unary predicate
``is a father'' abbreviates \(\exists z\,F(d,z)\), whereas
``is your father'' fixes that second argument at the owner. Possession
does not identify those two arguments. All ranges used in these
countermodels are nonempty; none of the failures rests on empty-term
conventions.
\end{proof}
\end{annexitem}

\needspace{8\baselineskip}
\begin{annexitem}{annex:p13.supposition.complements}{thm:p13-supposition-complements}{Obversion}
\begin{proof}
Fix a domain \(U\), and write \(\overline P=U\setminus P\).
If \(A(S,P)\), then \(S\subseteq P\), so
\(S\cap\overline P=\varnothing\), which is
\(E(S,\overline P)\). Replacing \(P\) by \(\overline P\)
and using \(\overline{\overline P}=P\) proves the second implication.
If \(I(S,P)\), choose \(s\in S\cap P\). This element is outside
\(\overline P\), so \(S\nsubseteq\overline P\), whence
\(A(S,\overline P)\) is false and \(O(S,\overline P)\) is true.
Again replace \(P\) by \(\overline P\) for the fourth implication.

For every converse use \(U=\{u\}\) and \(S=\varnothing\), with
either choice of \(P\). Both negative clauses E and O are true,
and both affirmative clauses A and I are false. This single construction
therefore refutes all four converses without changing the declared clauses.

Contraposition of a categorical interchanges its terms and complements
both. For the A-form take \(S=P=U=\{u\}\). Then \(A(S,P)\)
is true but \(A(\overline P,\overline S)=A(\varnothing,\varnothing)\)
is false. For the O-form take \(S=P=\varnothing\) in the same
nonempty \(U\). Then \(O(S,P)\) is true, while
\(O(\overline P,\overline S)=O(U,U)\) is false. These are the two
forms for which the source discusses conversion by contraposition.
If the operation is proposed without restriction for E or I as well,
the same two examples refute it: E holds at
\((\varnothing,\varnothing)\) but fails at \((U,U)\), and I holds
at \((U,U)\) but fails at \((\varnothing,\varnothing)\).
\end{proof}
\end{annexitem}


\needspace{8\baselineskip}
\begin{annexitem}{annex:p13.supposition.time}{thm:p13-supposition-time}{Two temporal countermodels}
\begin{proof}
Take one individual \(u\) existing at each of three times
\(-2,-1,0\), with \(0\) present. Let it be white only at \(0\)
and black only at \(-1\). The range of things that are now or formerly
white contains \(u\), and the range of formerly black things also
contains \(u\). Their intersection is nonempty, so the declared
ampliated reading of \emph{A white was black} is true. At no time is
\(u\) both white and black. Consequently the present-tense form was
never true. This refutes the rejected past-truth paraphrase without
using a failure of existence.

For the second example retain the same individuals and times, but let
\(u\) be a bishop only at \(-2\) and grey only at \(-1\).
The now-or-formerly-bishop range and the formerly-grey range each
contain \(u\). Their intersection again makes the ampliated sentence
true. The two properties never hold simultaneously. The construction
therefore distinguishes separate existential choices of past times
from one existentially chosen past time shared by both predicates.
It does not claim that either construal is compulsory in ordinary
speech. It checks the exact independent-range account under discussion
\cite[pp.~206--210]{Parsons2007FordhamSupposition}.
\end{proof}
\end{annexitem}

\needspace{8\baselineskip}
\begin{annexitem}{annex:p13.supposition.dependent-range}{thm:p13-supposition-dependent-range}{The owner must remain bound}
\begin{proof}
Let the farmers be \(f_1,f_2\), the donkeys \(d_1,d_2\), and
ownership between these ranges consist of the pairs
\((f_1,d_1),(f_2,d_2)\). Let only \(d_1\) run. Farmer \(f_1\)
has a nonempty donkey range and every member of that range runs.
Nevertheless \(d_2\) is a nonrunning donkey owned by a farmer.
Thus neither the range of all donkeys nor the range of donkeys owned
by some farmer can replace the particular witnessing farmer's range.
The anaphoric restriction cannot be discarded.

For the superior/inferior failure let the farmers again be \(f_1,f_2\),
and let the animals be \(h,d\). Only \(d\) is a donkey, and only
\(h\) runs. Let \(f_1\) own exactly \(h\) and \(f_2\) own
exactly \(d\). Every donkey is an animal, and the donkey range is
nonempty. Some farmer's every animal runs, witnessed by \(f_1\).
But no farmer witnesses the corresponding non-vacuous assertion about
his donkeys: \(f_1\) has none, while \(f_2\)'s donkey does not run.
The general inclusion of donkey in animal does not provide nonemptiness
of the inferior range for the same owner. This proves the stated
failure with exactly the affirmative existence convention of the body
\cite[pp.~255--257]{Parsons2007FordhamSupposition}.
\end{proof}
\end{annexitem}

\needspace{8\baselineskip}
\begin{annexitem}{annex:p13.supposition.four-modes}{thm:p13-supposition-four-modes}{Finite instances and negation duality}
\begin{proof}
Assume finite nonempty \(S\) and \(P\). Write the corresponding
listed members as \(s_i\) and \(p_j\). The listing is exhaustive;
all instances retain the reference information that these elements
belong to their ranges. To test failures of a formal inference, the
two term ranges vary independently and can have two or more members.
This avoids assigning a mode through an accidental identity between
two occurrences of the same word.

For the subject of A, an instance is \(s_i\in P\). All such
instances together are equivalent to \(S\subseteq P\), and
nonemptiness of \(S\) supplies the remaining A-conjunct. Thus A is
equivalent to their conjunction. Ascent from one instance is invalid:
take distinct \(a,b\), \(S=\{a,b\}\), \(P=\{a\}\).
The instance at \(a\) is true while A is false. This is W.

For the predicate of A, an instance is the assertion that every member
of \(S\) is the particular \(p_j\). Each instance entails A.
But take \(S=P=\{a,b\}\). A is true, while neither of those
uniform identity assertions is true. No disjunction of the instances,
and therefore no conjunction, follows from A. This is MC.

For I the subject instances are \(s_i\in P\). Their disjunction
is exactly \(S\cap P\ne\varnothing\). Conjunctive descent is
invalid in the first two-element example. The predicate instances
\(p_j\in S\) give the same argument with the roles interchanged.
Both occurrences are Det.

For E the subject instances are \(s_i\notin P\), and the
predicate instances are \(S\cap\{p_j\}=\varnothing\).
Each complete conjunction is equivalent to \(S\cap P=\varnothing\).
Neither permits ascent from an arbitrary instance. For the subject
use \(S=\{a,b\},P=\{a\}\), selecting \(b\); for the
predicate interchange these ranges, again selecting \(b\) from its
two-element range. E is false in either model although the selected
instance is true. Thus both occurrences are W.

For the subject of O, nonemptiness of \(S\) makes O equivalent to
\(S\nsubseteq P\). Its instances \(s_i\notin P\) therefore
have disjunction equivalent to O. In the first two-element example
one instance is true and another false, so conjunctive descent is
invalid. This is Det. Notice where nonemptiness was used: it would
be incorrect to extract this witness when \(S\) is empty.

For the predicate of O the instance at \(p_j\) is that some member
of \(S\) differs from \(p_j\). A witness in \(S\setminus P\)
supplies every such instance, so O descends conjunctively. Return
fails even from the conjunction. Take \(S=P=\{a,b\}\). For each
\(p_j\) the other element is a member of \(S\) different from it;
all instances hold. But \(S\subseteq P\) and \(S\ne\varnothing\),
so O is false. This is N and also exhibits the missing single common
witness behind the failed return.

Finally let \(F_i\) be any finite nonempty family of instances for
\(F\), with \(C=\bigwedge_iF_i\) and
\(V=\bigvee_iF_i\). The instances for the simultaneous negation
comparison are \(\neg F_i\). Their conjunction is \(\neg V\)
and their disjunction is \(\neg C\), by De Morgan's laws. Thus
\(F\Leftrightarrow V\) is equivalent to
\(\neg F\Leftrightarrow\neg V\), exchanging the Det return
pattern with the W return pattern. The failure of \(F\Rightarrow C\)
is, by contraposition, the failure of
\(\neg C\Rightarrow\neg F\), so the accompanying exclusion
of single-instance ascent is preserved as well. Similarly,
\(V\Rightarrow F\) together with failure of
\(F\Rightarrow V\) is equivalent to
\(\neg F\Rightarrow\neg V\) together with failure of
\(\neg V\Rightarrow\neg F\). This exchanges MC and N.
The argument is about these specified instances, not about arbitrary
opaque contexts in which negating the whole sentence would change
what counts as an instance.
\end{proof}
\end{annexitem}

\needspace{8\baselineskip}
\begin{annexitem}{annex:p13.supposition.prefix-rules}{thm:p13-supposition-prefix-rules}{Prefix rules}
\begin{proof}
Work in the declared fragment with fixed nonempty ranges and an
extensional terminal clause. A restricted existential is ordinary
existential quantification over its range, and a restricted universal
is ordinary universal quantification over that nonempty range. The
negative quantifier is the negation of the restricted existential.
Consequently its normalization to a restricted universal followed by
negation is valid. The dualities moving \emph{not} past \emph{some}
or \emph{every} and double-negation cancellation are likewise valid.
No rule permutes the ordered term occurrences or changes a range.

The normalization terminates: each passage of a negation through a
phrase reduces the number of phrases still to its right within that
negation's scope; cancellation removes signs. Equivalently, normalize
recursively from the innermost suffix, recording the parity of negations
passed by each phrase. Its final sign is its original affirmative sign
dualized at each such passage, and the terminal clause retains the
remaining parity. This specifies a unique designated result with the
original term order. It asserts no uniqueness among all logically
equivalent prefixes.

Classify a normalized occurrence by its quantifier sign and by whether
there is an earlier opposite sign. We verify the four outward operations
on this pair. Prepending an existential makes the new occurrence wide
existential. For an old occurrence it preserves both existential
classes; it changes wide universal to narrow universal, and an already
narrow universal remains narrow. These are the \emph{some} row.
Prepending a universal similarly changes wide existential to narrow
existential, preserves the other classes, and makes its own occurrence
wide universal. These are the \emph{every} row.

Prepending negation dualizes all signs in the suffix. Being opposite
is invariant under simultaneous dualization, so wide/narrow status is
unchanged while existential/universal status is reversed. Hence
\(\mathrm{Det}\leftrightarrow\mathrm W\) and
\(\mathrm{MC}\leftrightarrow\mathrm N\), exactly the \emph{not}
row. Finally \emph{no} is a new universal with negation before the
old suffix. Apply the just-proved negation row, then the universal row.
In the table's column order, the resulting old classes are respectively
\(\mathrm W,\mathrm N,\mathrm{MC},\mathrm{MC}\), as displayed.
The new term is wide universal.

Starting at the terminal clause and adjoining each sign proves the
claim by structural induction on the prefix. This is a complete
verification of the transition algorithm against the specified prefix
classification. The historical comparison is Parsons's refined
discussion at pp.~269--273; the proof does not extend its instance
tests to dependent ranges merely because the same string algorithm
can label their quantifiers.
\end{proof}
\end{annexitem}

\needspace{8\baselineskip}
\begin{annexitem}{annex:p13.supposition.relative-check}{thm:p13-supposition-relative-check}{Instances}
\begin{proof}
Let \(M,H,D\) be finite nonempty ranges of men, horses, and donkeys.
Let \(T(m,h)\) mean that the man sees the horse and \(R(d,h)\)
that the donkey sees it. The first sentence is
\[
 \exists m\in M\ \exists h\in H\ \exists d\in D\,
       (T(m,h)\land R(d,h)).
\]
For a fixed listed donkey \(d_i\), its instance leaves just the
first two existential quantifiers and uses \(R(d_i,h)\). The
original is equivalent to the disjunction of all these instances,
because the innermost existential ranges over exactly the listed
donkeys. Conjunctive descent fails: let one man see one horse and
let just one of two donkeys see it. Thus the pattern is Det.

The second sentence changes only the first existential to a universal:
\[
 \forall m\in M\ \exists h\in H\ \exists d\in D\,
       (T(m,h)\land R(d,h)).
\]
Any one-donkey instance implies the original, but the original need
not imply their disjunction. Take two members in each range, and let
\(T\) contain just \((m_1,h_1),(m_2,h_2)\), while \(R\)
contains just \((d_1,h_1),(d_2,h_2)\). The original is true.
The instance at \(d_1\) fails for \(m_2\), and the instance
at \(d_2\) fails for \(m_1\). This is MC.

For the closing example let \(H\) be the horse range, let \(A\)
be the animal range, and let \(R(h,d)\) be seeing. Its explicitly
extensional reading is
\[
 F=\exists h\in H\,
       \bigl(h\in A\land\forall d\in D\ R(h,d)\bigr).
\]
The instance at \(d_i\) is
\(F_i=\exists h\in H\,(h\in A\land R(h,d_i))\).
The one horse witnessing \(F\) witnesses every \(F_i\), so
\(F\Rightarrow\bigwedge_iF_i\). Return fails with two horses
and two donkeys: set \(A=H=\{h_1,h_2\}\),
\(D=\{d_1,d_2\}\), and let seeing contain just
\((h_1,d_1),(h_2,d_2)\). Each instance is true but no horse
sees every donkey, so \(F\) is false. Hence the occurrence is
narrowly distributed. In particular it cannot be merely confused,
which would prohibit precisely this valid conjunctive descent.

The held manuscript prints this sentence with a merely confused
classification at p.~276. The displayed countermodel and general
descent identify the correction for the declared relational reading.
No collation with the distinct final published manifestation, and no
claim to repair its entire non-main-term theory, is implicit in this
authorial check.
\end{proof}
\end{annexitem}
